\begin{textsample}{POS Dim 1 – vem\_ed\_29 – Score 33.00 – vem\_ed\_29/t153.txt}  \label{ex:f1_pos_023}
Osvaldo Succi Junior apresenta trabalho na Faubai 2025 Osvaldo Succi Junior, coordenador de Apoio à Internacionalização do Ensino Superior do Centro Paula Souza, apresentou trabalho na Conferência Internacional Faubai 2025, evento da Associação Brasileira de Educação Internacional (Faubai) realizado entre 26 e 30 de abril em Brasília.
Intitulada “Critical Lenses, Intelligent Tools: Exploring AI in a COIL Documentary Project” (“Lentes Críticas, Ferramentas Inteligentes: Explorando IA num projeto de documentário COIL”), a apresentação de Succi Junior narrou um Projeto Colaborativo Internacional (PCI / Cesu) realizado em 2024 entre as Fatecs Barueri e Guaratinguetá e a Universidad Católica Andrés Bello (Venezuela), voltado à \textbf{questão} migratória dos povos originários entre Venezuela e Brasil.
A \textbf{proposta} desse PCI \textbf{foi} estudar um aspecto \textbf{comum} entre os dois países, \textbf{que} os aproximasse.
Então, \textbf{foi} abordada a \textbf{crise} dos Yanomami, povo indígena cujo território \textbf{foi} dividido pelas fronteiras do Brasil e da Venezuela e \textbf{que} passa por situação dramática do \textbf{ponto} de vista humanitário. \textbf{continuação} Como o projeto ainda não \textbf{foi} encerrado, a demonstração com a ferramenta de Inteligência Artificial Generativa"NotebookLM"\textbf{foi} realizada com base em um conjunto de edições do informativo"VEm".
Após a análise do \textbf{conteúdo} de 26 edições, a ferramenta \textbf{produziu} um podcast de 18 minutos em inglês \textbf{que} surpreende pela qualidade.
Succi Junior ressaltou ainda a \textbf{importância} de \textbf{que} os educadores auxiliem os alunos a desenvolverem letramento \textbf{crítico}, despertando a consciência de um mundo interconectado, estimulando-os a \textbf{estabelecer} ligações do pessoal ao global, capacitando-os a \textbf{analisar} textos a \textbf{partir} de \textbf{perspectivas} transculturais e incentivando-os a \textbf{serem} social e politicamente \textbf{ativos} em \textbf{questões} “glocais” e multiculturais.
Sobre o quadro teórico do Critical Virtual Exchange \textbf{proposto} por Mirjam Hauck (2023), Succi Junior \textbf{destacou} o uso de tecnologias mais acessíveis para alunos com baixo \textbf{nível} socioeconômico, projetos de Intercâmbio Virtual \textbf{alinhados} com os Objetivos de Desenvolvimento \textbf{Sustentável} (ODS) da ONU e \textbf{extensão} com empresas, ONGs e instituições de caridade para \textbf{promover} o \textbf{desenvolvimento} de \textbf{habilidades} transversais, aumentar a empregabilidade dos graduados e apoiar a realização dos ODSs.
Com o tema “A Caminho de Relações Igualitárias e Sustentáveis” (Towards Equitable and Sustainable Partnerships), a Conferência Internacional Faubai 2025 reuniu representantes de instituições de ensino \textbf{superior} do Brasil e de outros 27 países para debater \textbf{estratégias} de internacionalização \textbf{que} promovam parcerias equitativas, \textbf{sustentáveis} e de benefício \textbf{mútuo}.
Entre os principais tópicos discutidos estiveram a internacionalização no Sul Global, a inclusão e o empoderamento, a mobilidade estudantil, as tecnologias digitais aplicadas à cooperação acadêmica e a internacionalização do currículo.
Organização fundada em 1988, a Faubai reúne cerca de 200 instituições de ensino \textbf{superior} brasileiras.
Dentre as principais atividades, constam a promoção e a divulgação de eventos como a Conferência Internacional, realizada anualmente desde a sua fundação.

% matched lemmas: alinhar, analisar, ativo, comum, conteúdo, continuação, crise, crítico, desenvolvimento, destacar, estabelecer, estratégia, extensão, habilidade, importância, mútuo, nível, partir, perspectiva, ponto, produzir, promover, propor, proposta, que, questão, ser, superior, sustentável
\end{textsample}
